\anexo{Entrevista 1: Vendedor de Mercado Libre (foco: la problemática)}

Como complemento cualitativo de la encuesta, se realizaron tres entrevistas semiestructuradas a usuarios reales de Mercado Libre: dos vendedores y una compradora, seleccionados por muestreo intencional entre personas con actividad verificable en la plataforma (no se buscan perfiles profesionales ni expertos, sino usuarios comunes, en línea con el público objetivo de la herramienta). El propósito es profundizar en el porqué de los patrones cuantitativos observados en la encuesta: la detección tardía de tendencias (88,9\% combinado), la alta disposición a usar una herramienta automatizada (68,9\%) y el peso de las señales sociales en la decisión de compra (47,5\%). Cada entrevista se registró con consentimiento informado y se analiza mediante síntesis temática en la Capítulo~\ref{chapter06} (User Research). Los nombres utilizados en las tres transcripciones (este anexo y los Anexos E y F) son pseudónimos elegidos para preservar el anonimato de los entrevistados (Sección~\ref{sec:viabilidad-legal}); no se registran apellidos, datos de contacto ni otra información identificatoria. A continuación se presenta el resumen completo de la primera entrevista, pregunta por pregunta, reportando en tercera persona lo que respondió el entrevistado.

\emph{Entrevistado}: Martín, 34 años. Vende en Mercado Libre hace 4 años, rubro bazar y hogar / gadgets ``virales''.

\emph{Objetivo}: entender el proceso real de decisión de stock y dimensionar el costo de detectar tendencias tarde.

\begin{enumerate}
	\item \textbf{¿Hace cuánto vendés en Mercado Libre, en qué rubro y qué tan seguido incorporás productos nuevos?} Contó que vende hace 4 años en el rubro bazar y hogar, principalmente gadgets de cocina y productos que se vuelven virales (freidoras, organizadores, luces). Suma productos nuevos una o dos veces por mes, cuando algo se repite mucho en redes sociales.
	\item \textbf{Contame la última vez que sumaste un producto nuevo: ¿cómo decidiste que convenía?} Relató que la última incorporación fue una organizadora de placard que vio explotar en TikTok: la vio en más de diez videos distintos en una semana, buscó en Mercado Libre qué oferta había, notó que había pocos vendedores y que vendían poco (menos de 100 unidades), pidió precio a su importador y calculó el margen con envío. Aclaró que no usó ningún sistema para decidir, sino intuición combinada con lo que veía en redes.
	\item \textbf{¿Qué fuentes usás hoy para saber qué se está vendiendo o qué está creciendo?} Mencionó TikTok e Instagram (reels de ``cosas de hogar''), Google Trends, los ``Buscados'' de Mercado Libre, seguir a cuatro o cinco vendedores grandes para ver qué publican, y un grupo de WhatsApp de vendedores donde se intercambian tips. Señaló que Mercado Libre Analytics solo le sirve para ver su propia venta, no la del mercado en general, y que eso es justamente lo que le falta.
	\item \textbf{¿Cuánto tiempo por semana le dedicás a investigar tendencias o productos para stockear?} Estimó unas 3 o 4 horas por semana, pero las considera mal invertidas: scrollea redes, mira rankings y compara precios a mano, sin llegar a una respuesta clara al final del proceso.
	\item \textbf{¿Te pasó de llegar tarde a una tendencia? Contame el caso y qué consecuencia tuvo.} Sí: contó el caso de las luces proyector de estrellas (galaxy light). Las vio en TikTok en marzo, dudó una semana y pidió muestras; para cuando tuvo stock disponible, ya había 200 publicaciones y el precio de referencia había bajado un 30\%. Compró 40 unidades, vendió 12 a buen precio y liquidó el resto casi a costo. Reconoció que perdió plata y que quedó con miedo de meterse en la siguiente tendencia, un miedo que también tuvo costo, porque a las últimas dos tendencias que vio directamente no entró.
	\item \textbf{Si esa información la hubieras tenido dos semanas antes, ¿qué habrías hecho distinto?} Respondió que habría comprado el doble de stock al precio viejo, publicado antes de que la categoría se saturara y cobrado entre 2 y 3 veces más por unidad, además de arrancar posicionado mientras la demanda todavía estaba creciendo y no cuando ya había 200 competidores. Estimó la diferencia de rentabilidad en una relación de 3 a 1.
	\item \textbf{¿Cómo distinguís una moda pasajera de una tendencia que se va a sostener?} Explicó que mira tres cosas: si la tendencia aparece en cuentas grandes y no solo en cuentas chicas, si se sostiene más de 3 o 4 semanas en redes y en Google Trends (y no es solo un pico de un día), y si en el marketplace todavía hay pocos vendedores pero la búsqueda ya viene creciendo. Dijo que desconfía cuando solo encuentra videos de cuentas chicas con ``link en bio''.
	\item \textbf{¿Qué información te falta hoy y no encontrás en ningún lado?} Señaló que le falta poder ver la velocidad a la que crece la demanda antes de que se note en las ventas: según su experiencia, no existe ningún lugar que indique que un producto está acelerando en tiempo real. Mercado Libre muestra rankings, pero sin velocidad ni historial ni precios históricos, y las ventas de la competencia que se muestran son una estimación que decide la plataforma, no un dato real verificable.
	\item \textbf{¿Probaste alguna herramienta o método para esto (planillas, apps, seguir cuentas)? ¿Qué te aportó?} Contó que probó planillas propias con datos de Google Trends, pero las abandonó a las dos semanas porque daban mucho trabajo; también probó seguir cuentas, el ranking de Mercado Libre y el grupo de WhatsApp de vendedores. De todo eso, lo que más le aportó fue el grupo de WhatsApp, porque hay gente que ya probó cosas y comparte qué funciona y qué no. Consideró que TikTok le sirve para descubrir productos, pero le hace perder mucho tiempo.
	\item \textbf{¿Qué tendría que tener una herramienta para que la uses todas las semanas?} Mencionó tres condiciones: que la herramienta avise sola cuando algo empieza a acelerar (porque si tiene que entrar a revisarla, no la va a usar), que muestre el crecimiento en números simples y estime cuánto puede ganar (margen frente a saturación), y que le insuma menos de 10 minutos por semana. Advirtió que si tuviera que cargar datos manualmente, la abandonaría, como pasó con las planillas.
\end{enumerate}
